% doc/metrics-protocol.tex -- appendix documenting the literature search
% behind the claims made in the `learnlog metrics' chapter.  Plain LaTeX:
% no preamble, no \documentclass; \input after \appendix in doc/learnlog.nw.
% The per-reference provenance blocks summarised here live in
% doc/bibliography.bib, above each entry.

% A local helper, used nowhere else in the document: verbatim queries are
% set in the typewriter font, whose interword space has no stretch, so a
% long query cannot be justified and overflows the measure.  \qy gives such
% a query a stretchable interword space and changes nothing else.
\providecommand{\qy}[1]{{\ttfamily\spaceskip=0.5em plus 0.5em minus 0.17em
  \hyphenpenalty=10000 #1}}

% Long verbatim URLs below are recorded exactly as they were fetched, and
% several of them contain runs of hyphenated words that url.sty will not
% break by default.  Allow breaks after hyphens for the duration of this
% appendix; it is restored at the end of the file.
\makeatletter
\def\do@url@hyp{\do\-}
\makeatother
\Urlmuskip=0mu plus 1mu\relax

\chapter{Literature search protocol for the metrics}
\label{app:metrics-protocol}

The metrics that \qy{learnlog metrics} computes are not learnlog's
inventions.  Each one is defined in a published paper, each carries
parameters that a reimplementation must choose, and each has a literature
saying what it does and does not measure.  This appendix records how that
literature was found, what was verified, and what was not, so that a reader
can audit every backed claim rather than take it on trust.

\section{The protocol}
\label{app:protocol:method}

One protocol governs every claim in this appendix, and it is stated here
once rather than repeated in each chapter.

\begin{description}
\item[Find, in both directions] Every claim was searched twice: once for
  literature that supports it, and once, with a separate
  \emph{counter-search}, for literature that refutes or qualifies it.  A
  counter-search that returned nothing is recorded as an empty
  counter-search, not omitted.
\item[Verify at the source] No claim rests on a title or a search-result
  snippet.  Where the full text could be obtained it was read; where only
  the publisher-supplied abstract could be obtained, the claim is stated at
  abstract strength and says so; where neither could be obtained, the source
  is listed in \cref{app:protocol:dropped} and backs nothing.
\item[Record] Each reference carries a provenance block in
  \nolinkurl{doc/bibliography.bib} immediately above its entry, giving the
  claim it backs, the verbatim query that found it, why it was picked over
  the alternatives, the verbatim supporting passage, what was verified, and
  the counter-search outcome.
\end{description}

The chapters that follow are organised by \emph{claim}, not by search
episode: one chapter per claim, each stating the claim as a research
question and answering it from the evidence gathered for that claim alone.

\section{Sessions, date and providers}
\label{app:protocol:sessions}

All searching was done on \textbf{2026-08-31} under three named
\texttt{scholar} sessions, whose exported records are the audit trail behind
this appendix:

\begin{description}
\item[\texttt{learnlog-metrics-eq}] Error Quotient, Repeated Error Density,
  the Watwin score and the Normalized Programming State Model; 24 recorded
  searches.
\item[\texttt{learnlog-metrics-time}] time-on-task estimation, session
  segmentation and effort proxies; 10 recorded searches over 6\,505
  candidate records.
\item[\texttt{learnlog-metrics-code}] static code metrics, the
  maintainability index and linters for novice code; 11 recorded searches
  over 139 screened papers.
\end{description}

Each session is stored as a JSON file named after it, under
\nolinkurl{~/.local/share/scholar/review_sessions/}.

Provider health was checked with \qy{scholar providers check} before
searching, and the result materially limits what a negative finding in this
appendix means:

\begin{description}
\item[\texttt{s2} --- dead in all three sessions] The Semantic Scholar key
  was rejected with HTTP~403 (\qy{key rejected --- HTTP 403 --- check or
  replace S2\_API\_KEY}).  Semantic Scholar is the strongest provider for
  the ACM proceedings that dominate this literature, so its absence is the
  single largest gap in coverage.  It was partly compensated by direct
  \nolinkurl{api.semanticscholar.org} retrieval outside the session.
\item[\texttt{openalex} --- reported healthy, effectively absent] It
  returned HTTP~429 or 503 on every \texttt{scholar} query in the EQ and
  time sessions, despite passing the health check; it behaved normally in
  the code session.  Direct REST retrieval
  (\nolinkurl{api.openalex.org/works/doi:} followed by the DOI) worked
  throughout
  and was used instead.
\item[\texttt{dblp} --- degraded] Intermittent HTTP~500 and 503, and a hard
  limit of about three terms per query (every term must match, implicit
  AND), which forced several precise queries to be shortened.  Its keyword
  index also conflates \enquote{error quotient} with numerical-analysis
  \enquote{difference quotients}.
\item[\texttt{wos} --- healthy, and it carried the sessions] The field-tagged
  \texttt{TS=(\ldots)} topic syntax was the only reliable discovery route for
  counter-evidence in all three sessions.
\item[\texttt{ieee} --- healthy] One transient HTTP~502.
\item[\texttt{scopus} --- healthy but useless here] Every query in the EQ
  session returned recency-ranked, off-topic records (2026 nursing, finance
  and materials papers); all were discarded.
\end{description}

\paragraph{A method limitation, stated plainly.}
Natural-language counter-queries did not work on these providers.  In the
code session, \qy{cyclomatic complexity critique validity} returned zero
hits and \qy{maintainability index criticism limitations} returned only
off-topic records (soil erosion, the Barthel Index, positive psychology),
while the field-tagged Web~of~Science queries that followed surfaced the
whole critical literature.  Anyone reproducing this work should phrase
counter-searches as \texttt{TS=(\ldots)} and should not read a
natural-language \enquote{nothing found} as a negative result.  Two of the
largest databases were silently absent from the sessions, so a
\enquote{not found} recorded here is weak evidence; every load-bearing item
was therefore additionally retrieved by DOI through Crossref or OpenAlex.

\section{Complements outside \texttt{scholar}}
\label{app:protocol:complements}

The protocol allows any discovery route provided the route is recorded.
The complements below are every non-\texttt{scholar} tool invocation that
contributed to the bibliography, grouped by session; each gives the tool,
the verbatim query or URL, and the outcome.  All are dated 2026-08-31.

\subsection{Session \texttt{learnlog-metrics-eq}}

\begin{description}
\item[WebSearch] Fifteen queries, verbatim:
  \begin{itemize}
  \item \qy{Jadud "Methods and tools for exploring novice compilation
    behaviour" ICER 2006 error quotient} --- located the author-hosted
    thesis; retained.
  \item \qy{Watson Li Godwin "Predicting performance in an introductory
    programming course\ldots" Watwin ICALT 2013} --- retained.
  \item \qy{Becker "A new metric to quantify repeated compiler errors
    for novice programmers" ITiCSE 2016 repeated error density} ---
    surfaced the RPTEL review and ProgSnap2Analysis; retained.
  \item \qy{Carter Hundhausen Adesope "normalized programming state
    model" ICER 2015} --- retained.
  \item \qy{Watson "Watwin" algorithm score "time penalty" quartile
    compilation pairs BlueFix thesis pdf} --- located the Durham e-thesis;
    retained.
  \item \qy{Petersen Spacco Vihavainen "An exploration of error quotient
    in multiple contexts" Koli Calling 2015 pdf} --- no open copy.
  \item \qy{Ihantola "Educational data mining and learning analytics in
    programming\ldots" ITiCSE-WGR 2015 DOI} --- retained.
  \item \qy{Jadud Dorn "Aggregate compilation behavior: findings and
    implications from 27,698 users" ICER 2015 DOI} --- retained; also
    surfaced the DELFI Python paper.
  \item \qy{Becker RED "repeated error density" formula "Rn" novice
    compiler errors pdf brettbecker} --- no open copy.
  \item \qy{"normalized programming state model" NPSM states
    "syntactically" "semantically" correct incorrect programming states
    seconds} --- retained, for the state descriptions.
  \item \qy{Karvelas Dillane Becker "Compile Much" UKICER 2020 novice
    compilation error quotient pdf ucd repository} --- abstract only.
  \item \qy{Karvelas Li Becker "compilation mechanisms" SIGCSE 2020
    "error quotient" automatic manual compilation results} --- abstract
    only.
  \item \qy{Petersen Spacco Vihavainen error quotient multiple contexts
    findings "free parameters" conclusion EQ did not predict} --- retained;
    led to Ahadi et al.\ 2018.
  \item \qy{"error quotient" penalty "13" "3" normalize 16 Petersen 2015
    many small assignments adapted penalties} --- \textbf{empty}; the
    targeted search for the penalty values found nothing.
  \item \qy{Spacco "error quotient" Koli 2015 Petersen Vihavainen author
    copy pdf knox.edu OR cs.helsinki.fi OR utoronto} --- reached Petersen's
    publications page; no PDF.
  \end{itemize}
\item[alphaXiv] \texttt{discover\_papers} with keywords \qy{["error
  quotient", "Jadud", "compilation behaviour", "novice programmers",
  "Watwin", "repeated error density", "Python"]} and the question
  \qy{"Error Quotient penalty parameters for novice programming, and
  adaptations of EQ to Python or many-small-assignment contexts"} --- did
  not hold Petersen et al.\ 2015 (not an arXiv paper); surfaced PyMETA,
  retained only as a pointer.  Then \texttt{get\_paper\_content} on
  \url{https://arxiv.org/abs/2606.30610} --- retained.
\item[WebFetch] \url{https://jadud.com/dl/pdf/jadud-dissertation.pdf};
  \url{https://etheses.durham.ac.uk/11141/};
  \url{https://d-nb.info/1212501446/34};
  \url{https://arxiv.org/pdf/2404.05988};
  \url{https://dl.gi.de/handle/20.500.12116/44524};
  \url{https://kth.diva-portal.org/} record \texttt{diva2:1877365};
  \url{https://utmandrew.bitbucket.io/pubs.html};
  \url{https://onlinelibrary.wiley.com/doi/full/10.1002/cae.22640}; and the
  author-hosted copy of Hundhausen et al.\ 2017 --- all retained.
\item[REST retrieval] Each of these takes the item's DOI as its last path
  element: \nolinkurl{api.crossref.org/works/} for 27 DOIs;
  \nolinkurl{api.openalex.org/works/doi:} for 7;
  \nolinkurl{api.datacite.org/dois/10.18420/delfi2024_40};
  \nolinkurl{api.unpaywall.org/v2/} for 5;
  and \nolinkurl{api.semanticscholar.org/graph/v1/paper/DOI:} ---
  retained; this is where most abstract-level verification came from once
  \texttt{s2} and \texttt{openalex} were unavailable inside the session.
\item[GitHub raw] The ProgSnap2Analysis reference implementations
  \texttt{eq.py}, \texttt{red.py} and \texttt{watwin.py} --- retained; they
  are the only source for several penalty values, and their own comments
  document which parts are guesses.
\item[Local first-party sources] The user's own manuscript sources
  \nolinkurl{edit-run-error-repeat/conference_101719.tex} and
  \texttt{base.bib} --- retained; the two-parameter EQ penalties are
  verified verbatim there and nowhere else.
\item[Failed tool invocations] \qy{scholar enrich learnlog-metrics-eq}
  timed out at 300\,s and was abandoned; \qy{scholar pdf quote} on the
  ACM PDF of Karvelas et al.\ 2020 returned HTTP~403 with no passages.
\end{description}

\subsection{Session \texttt{learnlog-metrics-time}}

\begin{description}
\item[WebSearch] Twelve queries, verbatim:
  \begin{itemize}
  \item \qy{Leinonen Castro Hellas "time-on-task" metrics predicting
    performance programming SIGCSE} --- retained; corrected a venue error
    (ICER 2017, not ICER 2021).
  \item \qy{Kovanovi\'c "Penetrating the black box of time-on-task
    estimation" LAK 2015} --- retained; led to the Edinburgh open-access
    PDF.
  \item \qy{Leinonen Hellas "comparison of time metrics in programming"
    ICER 2021} --- retained; surfaced the EDM 2021 paper.
  \item \qy{Hart Warren Edwards "Accurate Estimation of Time-on-Task
    While Programming" pdf digitalcommons} --- retained; surfaced the PLOS
    ONE 2025 follow-up.
  \item \qy{ProgSnap2 specification SessionID "Session" event definition
    documentation} --- \textbf{distrusted}: it asserted the 20-minute rule
    with no primary source.
  \item \qy{progsnap2 specification pdf "SessionID" column optional
    "Session.Start"} --- distrusted, same reason.
  \item \qy{"ProgSnap2" "SessionID" "20 minute" gap between consecutive
    events session} --- distrusted, same reason; all three were resolved by
    reading Price et al.\ directly instead.
  \item \qy{Edwards Snyder Perez-Quinones "Comparing effective and
    ineffective behaviors of student programmers" ICER 2009 doi} ---
    retained; corrected a widely circulated wrong DOI.
  \item \qy{Petersen Spacco Vihavainen "An Exploration of Error Quotient
    in Multiple Contexts" pdf session 20 minutes} --- \textbf{dropped}; no
    full text, so the threshold's origin stays a citation of a citation.
  \item \qy{Jadud "Methods and tools for exploring novice compilation
    behaviour" 2006 full text pdf error quotient definition} ---
    \textbf{dropped}; no full text obtained.
  \item \qy{Snider B\"alter Bosk "Edit, Run, Error, Repeat"
    time-on-task error quotient Koli Calling 2023 abstract} ---
    \textbf{dropped}; no abstract retrievable from any provider.
  \item \qy{"Comparison of Time Metrics in Programming" Leinonen
    Lepp\"anen Ihantola Hellas pdf helda} --- retained; located the
    author's PDF.
  \end{itemize}
\item[WebFetch, metadata] Again by DOI:
  \nolinkurl{api.crossref.org/works/} for seventeen DOIs;
  \nolinkurl{api.openalex.org/works/doi:} for three
  (verbatim abstracts);
  \nolinkurl{api.semanticscholar.org/graph/v1/paper/DOI:} for four
  (two abstracts obtained, two unavailable).
\item[WebFetch, full text]
  \url{https://pure.ed.ac.uk/ws/files/19361040/lak15pen.pdf} (10 pp.\ read);
  \url{https://juholeinonen.com/assets/pdf/leinonen2022time.pdf} (4 pp.);
  \url{https://juholeinonen.com/assets/pdf/leinonen2017comparison.pdf}
  (4 pp.);
  \url{https://educationaldatamining.org/EDM2021/virtual/static/pdf/EDM21_paper_176.pdf}
  (4 pp.);
  \url{https://ayaankazerouni.org/papers/2020/iticse20-progsnap2.pdf}
  (7 pp.); \url{https://digitalcommons.usu.edu/computer_science_stures/24}
  (abstract);
  \url{https://journals.plos.org/plosone/article?id=10.1371%2Fjournal.pone.0330758}
  (open-access article page) --- all retained.
\item[WebFetch, dropped] \url{https://github.com/CSSPLICE/progsnap2} --- the
  repository README does not define \texttt{SessionID}.
\item[Direct API query]
  \nolinkurl{dblp.org/search/publ/api?q=Fine-Grained+Versus+%
  Coarse-Grained}\ldots{} --- retained; it confirmed that \emph{no}
  page numbers exist for the EDM 2021 paper, so none were invented.
\end{description}

\subsection{Session \texttt{learnlog-metrics-code}}

\begin{description}
\item[Crossref API] \nolinkurl{api.crossref.org/works/} for
  \nolinkurl{10.1109/TSE.1976.233837} (McCabe), \nolinkurl{10.1109/2.303623}
  (Coleman et al.), \nolinkurl{10.1049/sej.1988.0003} (Shepperd),
  \nolinkurl{10.1145/3059009.3059061} (Keuning et al.\ 2017),
  \nolinkurl{10.1109/ICSM.1992.242525} (Oman \& Hagemeister),
  \nolinkurl{10.1145/3231711} (Keuning et al.\ 2018),
  \nolinkurl{10.1002/smr.1760} (Landman et al., volume, issue and pages
  verified) and
  \nolinkurl{10.1109/TSE.1983.236460} (metadata only; content unverified);
  plus the title query
  \nolinkurl{works?query.bibliographic=Designing+a+rubric+}\allowbreak%
  \nolinkurl{for+feedback+on+code+}\allowbreak%
  \nolinkurl{quality+in+programming+courses}, which retained Stegeman et al.\
  2016 and surfaced Bobadilla et al.\ 2024.
\item[OpenAlex API] \texttt{works/doi:} for
  \nolinkurl{10.1145/3639474.3640051}, \nolinkurl{10.1145/3545945.3569829},
  \nolinkurl{10.1002/smr.1760}, \nolinkurl{10.1109/SDS59856.2023.10329116},
  \nolinkurl{10.1109/TSE.1983.236460},
  \nolinkurl{10.1109/FIE63693.2025.11328661},
  \nolinkurl{10.1145/3287324.3287503} and
  \nolinkurl{10.1145/2960310.2960325} ---
  abstracts retained; the last reported \qy{oa\_status: closed} and was
  therefore dropped as content-unverifiable.
\item[Unpaywall] \nolinkurl{api.unpaywall.org/v2/10.1145/3639474.3640051} ---
  only an ACM link (403) and a repository landing page; full text
  unobtainable.
\item[WebSearch] \qy{van den Aker Rahimi "Design principles for
  generating and presenting automated formative feedback on code quality
  using software metrics" pdf repository} --- retained, via the Open
  Universiteit research portal.  \qy{Shen Conte Dunsmore "Software
  Science Revisited" 1983 Purdue technical report pdf "software science"}
  --- located Purdue e-Pubs TR~81-376, which returns HTTP~403 to every
  fetch; content dropped.  \qy{Aivaloglou Hermans "How Kids Code and How
  We Know" Scratch cyclomatic complexity 78\% scripts pdf} --- no
  open-access primary; the figure is cited as reported by Keuning et al.
  \qy{SEI maintainability index formula "2.46" OR "2.4" perCM 171 - 5.2
  ln aveV Software Engineering Institute} --- superseded by reading Coleman
  et al.\ directly.  \qy{radon 6.0.1 Python 3.13 compatibility issue
  mando} --- background only.
\item[WebFetch] \url{https://research.ou.nl/en/publications/design-principles-for-generating-and-presenting-automated-formati}
  (verbatim abstract, retained);
  \url{https://learn.microsoft.com/en-us/visualstudio/code-quality/code-metrics-maintainability-index-range-and-meaning}
  (the Visual Studio variant of the maintainability index and its
  thresholds, retained);
  \url{https://openlibrary.org/search.json?q=Elements+of+Software+Science+Halstead}
  (publisher and ISBN for Halstead's monograph, retained);
  \url{http://www.ecs.csun.edu/~rlingard/comp589/ColemanPaper.pdf}
  (\textbf{full text read}, p.~46);
  \url{https://radon.readthedocs.io/en/latest/intro.html} and the PyPI JSON
  for \texttt{radon} and \texttt{mando} (tooling facts).
\item[WebSearch then \texttt{curl}]
  \url{http://www.literateprogramming.com/mccabe.pdf} ---
  \textbf{McCabe 1976 full text obtained and read} (13 pp.).  This is the
  source of the counting rule in \cref{app:claim:cc-definition}.
\item[GitHub raw and API] \texttt{rubik/radon@master}:
  \texttt{radon/visitors.py}, \texttt{radon/metrics.py},
  \texttt{radon/raw.py}, \texttt{radon/complexity.py},
  \texttt{docs/intro.rst}, \texttt{CHANGELOG}, \texttt{setup.py}; and
  \texttt{rubik/mando@v0.7.1/mando/core.py} --- read as the source of truth
  for what the tool actually counts, in preference to its documentation.
  \nolinkurl{api.github.com/repos/rubik/radon/commits} gave the last commit
  date.
\end{description}

\section{Sources dropped as content-unverified}
\label{app:protocol:dropped}

The following works exist and are bibliographically verified, but their
content could not be read.  None of them backs a claim in this appendix
beyond what is stated here; where one appears in
\nolinkurl{doc/bibliography.bib}, its provenance block says exactly which
strength it is cited at.

\begin{description}
\item[Petersen, Spacco \& Vihavainen 2015] The Koli Calling paper
  \autocite{petersen2015exploration}: full text unobtainable, ACM~403 to
  every route including the free-access \nolinkurl{?cid=} URL from the
  author's own publications page, Unpaywall
  reports it is not open access, four guessed author-copy URLs returned 404
  or 403, and alphaXiv does not hold it.  Only the publisher-supplied
  abstract was verified.  Consequently the \(13/3/{\div}16\) penalty
  attribution in \cref{app:claim:eq-parameters} is \emph{not} verified at
  source, and neither is the 20-minute value in
  \cref{app:claim:idle-threshold}.
\item[Watson, Li \& Godwin 2013] The ICALT paper
  \autocite{watson2013predicting}: full text 403 from every route.  The
  Watwin penalty values are therefore reimplementation-derived, not
  paper-verified.  Watwin is out of scope for the implementation for exactly
  this reason.
\item[Becker 2016] The ITiCSE paper \autocite{becker2016new}: ACM~403, and
  the green open-access copy at the issuing university returned HTTP~405.
  The RED
  formula is \emph{triangulated} across three independent sources rather
  than read in the original (\cref{app:claim:red-formula}).
\item[\enquote{Watwin Revisited}] Listed in Watson's thesis under work
  \textcquote{watson2015thesis}{in preparation for submission, or currently
  under review}.  No published version exists in
  Crossref.  It is not cited anywhere.
\item[Aivaloglou \& Hermans 2016] Closed access with no open copy anywhere
  (OpenAlex \qy{oa\_status: "closed"}).  The finding that most Scratch
  scripts have cyclomatic complexity one is therefore attributed as
  \emph{reported by} Keuning et al.\ \autocite{Keuning2017}, in whose full
  text the sentence appears verbatim.
\item[Shen, Conte \& Dunsmore 1983] The published abstract is 80 words and
  stops before any critical content; IEEE and the Purdue technical-report
  version both return HTTP~403.  It is not used to back any critique of
  Halstead's measures, only the attribution of software science to
  Halstead's monograph \autocite{Halstead1977}.
\item[Snider, B\"alter \& Bosk 2023] The Koli Calling paper
  \autocite{snider2023edit}: its own full text was not obtained (ACM~403;
  the institutional repository record carries no abstract).  Nothing is
  attributed to its
  findings; the two-parameter EQ penalties come from the first-party
  manuscript source instead.
\item[Brocker \& Schroeder 2024] The DELFI paper
  \autocite{brocker2024correlation}: existence, authorship and metric
  coverage verified through DataCite and the publisher's landing page; the
  full text was not obtained, so their
  operational split between syntax and runtime errors is unverified.
\item[Abstract- or metadata-strength only] Jadud 2006; Ihantola et al.\
  2015; Carter et al.\ 2015; Tabanao et al.\ 2011; Ahadi et al.\ 2018; the
  three Karvelas papers of 2020 and 2022; Kohn 2019; Smith \& Rixner 2019;
  McCall \& K\"olling 2019; Morrison et al.\ 2023; Stegeman et al.\ 2016;
  van den Aker \& Rahimi 2024; Prajapati \& Acu\~na 2025; Jay et al.\ 2009;
  Sj\o berg et al.\ 2012; and Oman \& Hagemeister 1992.  Each is cited in
  the chapter that uses it and carries a provenance block in the
  bibliography saying what was and was not read; they are listed together
  here, without citations, because a wall of sixteen references would say
  less than the one fact they share.  No claim in this appendix exceeds
  abstract or metadata strength for any of them.
\item[PyMETA] An unrefereed arXiv preprint \autocite{li2026pymeta} from a
  commercial author group.  Cited only as evidence that grounding a Python
  error taxonomy in the official exception hierarchy is an approach in use.
\item[Not re-found under a recorded query] Edwards, Kandru \& Rajagopal
  2017 and Birillo et al.\ 2022 (Hyperstyle) were verified in an earlier,
  pre-protocol pass but did not resurface under any session query, so they
  carry no provenance block and are not cited.
\item[Discarded wholesale] Every \texttt{scopus} result in the EQ session
  (recency-ranked, off-topic).  The \qy{scholar enrich} run on that
  session timed out at 300\,s and was abandoned.
\end{description}

\section{Corrections carried into the bibliography}
\label{app:protocol:corrections}

Verification turned up seven errors in circulating metadata and in the
secondary literature.  They are fixed in \nolinkurl{doc/bibliography.bib} and
recorded here so the fixes are not silently undone.

\begin{enumerate}
\item \enquote{Comparison of Time Metrics in Programming}
  \autocite{leinonen2017comparison} is ICER 2017, not ICER 2021.
\item The DOI \nolinkurl{10.1145/1584322.1584331}, widely circulated for
  Edwards et al.\ 2009, is wrong: it resolves to a different paper in the
  same proceedings.  The correct DOI is \nolinkurl{10.1145/1584322.1584325}
  \autocite{edwards2009comparing}.
\item Leinonen et al.\ 2022 \autocite{leinonen2022timeontask} has an
  \emph{ACM Inroads} reprint at \nolinkurl{10.1145/3534564}; the SIGCSE
  original is cited.
\item Kovanovi\'c et al.\ 2015 \autocite{kovanovic2015penetrating} is
  pp.~184--193; some secondary sources give pp.~64--68.
\item The EDM 2021 paper \autocite{leinonen2021finegrained} has neither a
  DOI nor recorded page numbers, confirmed against dblp; none were invented.
\item Landman et al.\ \autocite{Landman2016} appeared online first on
  2015-12-09 and in issue 28(7) in July 2016; it is cited as 2016, which is
  why OpenAlex's 2015 disagrees.  Crossref stores no issued year at all for
  Oman \& Hagemeister \autocite{OmanHagemeister1992}; the year 1992 comes
  from the proceedings title and the IEEE record.
\item Jadud's scoring figure is a four-term cascade, not the three-term
  decision tree that the secondary literature reproduces
  (\cref{app:claim:eq-parameters}); and the open-access review's sentence
  \enquote{A program is semantically correct if it produces a runtime
  exception} is an inverted typo that contradicts its own state table, so it
  was not relied on.
\end{enumerate}

\chapter{The 20-minute idle threshold}
\label{app:claim:idle-threshold}

\paragraph{Question.}
\qy{learnlog metrics} cuts a run history into sessions wherever the gap
between one run finishing and the next starting reaches a threshold, and it
defaults that threshold to 20 minutes.  Is 20 minutes the established
convention for \emph{coarse} programming events, and does it originate with
Petersen, Spacco and Vihavainen?

\paragraph{Method.}
Session \texttt{learnlog-metrics-time}, 2026-08-31, providers \texttt{dblp},
\texttt{wos}, \texttt{ieee}, \texttt{openalex}, \texttt{scopus}.  Support
queries: \qy{time-on-task estimation from programming process data};
\qy{session segmentation inactivity threshold student log data};
\qy{session inactivity threshold}; \qy{time-on-task programming}.
Counter-queries: \qy{session timeout threshold arbitrary sensitivity};
\qy{time-on-task estimation validity limitations}.  Three WebSearch
queries about the ProgSnap2 specification asserted the rule without a
primary source and were distrusted; the question was settled instead by
reading the format paper's own full text at
\url{https://ayaankazerouni.org/papers/2020/iticse20-progsnap2.pdf}.

\paragraph{Results.}
The claim is \emph{supported} by one primary source, read in full.  Price et
al.\ state the convention and apply it:
\blockcquote[sec.~3.3.1]{price2020progsnap2}{ProgSnap2 allows a dataset to
provide a SessionID column, and we used this value when it was present.
Otherwise, we defined a session as a series of events with no more than a 20
minute gap between any two consecutive events, as in [23].}
The same reading \emph{qualifies} the claim in a way the secondary sources
did not: \texttt{SessionID} is an \emph{optional} column of the format, so
the 20 minutes is an analysis convention, not part of ProgSnap2.

Nothing \emph{refutes} the convention.  What the counter-searches produced
instead is a family of competing thresholds at other event granularities:
five minutes for keystrokes \autocite{hart2023accurate}, about ten minutes
for keystrokes \autocite{leinonen2021finegrained,leinonen2022timeontask},
three minutes for edit events \autocite{leinonen2017comparison}, sixty
minutes for \enquote{work sessions}, and ten to sixty minutes across
web-usage-mining practice \autocite{kovanovic2015penetrating}.  The
threshold scales with event sparsity, which is precisely why a value must
not be transplanted between granularities.

The \emph{attribution} half of the question is weaker.  Reference~[23] in
the passage above is Petersen et al.\ \autocite{petersen2015exploration},
but that paper's full text could not be obtained
(\cref{app:protocol:dropped}), so the origin is verified only as a citation
of a citation.

\paragraph{Conclusion.}
Yes on the convention, with a caveat on the attribution.  Twenty minutes is
the documented convention for coarse programming events --- the granularity
of a program run --- and it is the only such convention the searches found;
learnlog therefore adopts it as the default.  But it is an analysis
convention rather than a specification, it is granularity-specific, and the
attribution to Petersen et al.\ is the format paper's, not one this work
independently checked.  The documentation says \enquote{following Price et
al., who follow Petersen et al.} and never attributes the number to Petersen
directly.  Because the choice of estimation strategy demonstrably moves
results \autocite{kovanovic2015penetrating}, the threshold is exposed as an
option and printed in every report.

\chapter{Coarse time-on-task and achievement}
\label{app:claim:time-not-predictive}

\paragraph{Question.}
Does coarse, log-derived time-on-task predict student achievement?  If it
does not, what should a period summary headline instead?

\paragraph{Method.}
Session \texttt{learnlog-metrics-time}, 2026-08-31, same five providers.
Support queries: \qy{time-on-task programming}; \qy{novice
programmer effort metrics compilation events code churn}.  Counter-queries:
\qy{time on task does not predict performance} (the single largest
contributor to the session, adding 2\,770 records); \qy{time-on-task
estimation validity limitations}.  Rather than screening the whole session,
two targeted screens were run over its stored decisions: a known-item
presence check, which confirmed that the topic-driven queries had genuinely
surfaced twelve of fifteen target works (the other three were retrieved by
DOI), and a counter-evidence screen filtering titles on time-related terms
conjoined with accuracy- and validity-related terms.  Three hits; two
verified and retained.

\paragraph{Results.}
Three independent studies \emph{support} the negative claim.  Leinonen et
al.\ measured both granularities on one cohort, so the comparison is not
confounded by dataset:
\blockcquote[sec.~4.1]{leinonen2022timeontask}{the coarse-grained
time-on-task does not statistically significantly correlate with the weekly
exercise points (p = 0.39), while the correlation between the fine-grained
time-on-task and weekly exercise points is statistically significant (p =
1.6e-25) albeit weak (r = 0.35).}
Edwards et al.\ reach the same place from a within-subjects design over
89\,879 submissions by 1\,101 students, which controls for student ability:
\blockcquote{edwards2009comparing}{When students received A/B scores, they
started earlier and finished earlier than when the same students received
C/D/F scores. They also wrote slightly more program code. They did not
appear to spend any more time on their work, however.}
An earlier study in the same line found total time uncorrelated with exam
scores while active time correlated weakly, and gave the positive half of
the answer --- what \emph{does} track effort:
\textcquote[sec.~4.1]{leinonen2017comparison}{active time correlates with
total count of source code events (r = .79), edit event count (r = .72), run
event count (r = .53) and focus change event count (r = .54)}, concluding
that \textcquote[sec.~2]{leinonen2017comparison}{the amount of changes to the
source code is a good metric for time}.

Two sources \emph{qualify} the claim further rather than refuting it.
Krieter validated log-derived time against parallel screen recordings ---
the only external ground truth found --- and reports that
\textcquote{krieter2022areyoustillthere}{common online time estimation
strategies do not represent the online time for our students accurately}.
Hart et al.\ find that
\textcquote{hart2025multiday}{self-report and objective measures provide
information on distinct, different aspects} of time-on-task, which supports
showing a log-derived number \emph{alongside}, not instead of, what a
student believes they spent.

The nearest thing to a \emph{refutation} is Hilliger et al., where
surfacing time-on-task did change behaviour \autocite{hilliger2021evaluating};
but its audience was teachers and its data was student self-report, so it
does not touch a claim scoped to student-facing, log-derived displays.  It
is retained as the honest positive counterweight.

\paragraph{Conclusion.}
No --- coarse time-on-task is not an achievement predictor, and two
independent designs agree.  So the period report headlines counts and
timing: active days, when the work started and how it was spread, runs,
sessions, failed runs, lines changed, and only then time.  Two caveats.
First, Leinonen et al.'s \enquote{coarse-grained} is a keystroke-derived
span from first keystroke to first submission, whereas learnlog's number is
run-to-run active time with idle gaps excluded; the transfer is by analogy
about \emph{granularity}, not an identity of measures.  Second, Krieter's
result is a standing qualification on every number the command prints, which
is why the label is \enquote{active time (run-to-run)} rather than
\enquote{time spent}.

\chapter{The two-parameter Error Quotient}
\label{app:claim:eq-parameters}

\paragraph{Question.}
learnlog scores an EQ pair by adding 13 when both runs end in an error and 3
more when the error type repeats, then dividing by 16.  Is this
parameterisation attested, and does it come, as the user's own prior work
states, from Petersen, Spacco and Vihavainen?

\paragraph{Method.}
Session \texttt{learnlog-metrics-eq}, 2026-08-31.  Discovery queries:
\qy{error quotient} on \texttt{dblp}; \qy{error quotient novice
compilation behaviour} on \texttt{dblp}, \texttt{wos}, \texttt{ieee} and
\texttt{openalex}.  Counter-query: \qy{TS=("error quotient" AND
(limitation OR criticism OR replication OR context))} on \texttt{wos}.  The
targeted search for the numbers themselves,
\qy{"error quotient" penalty "13" "3" normalize 16 Petersen 2015 many
small assignments adapted penalties}, returned \emph{nothing}.  Five further
attempts to reach the Petersen et al.\ full text all failed
(\cref{app:protocol:dropped}).  Jadud's original scoring figure was read by
rendering it to an image at 140\,dpi, because text extraction scrambles the
flowchart.

\paragraph{Results.}
The parameterisation is \emph{verified verbatim}, but in the user's own
first-party manuscript source rather than in a retrieved publication.  It
reads: \enquote{In this work we use the two parameter EQ model, with a
difference error penalty of 13 and a same error penalty of 3 as suggested in
previous research}, citing Petersen et al.  The bibliographic record of the
companion Koli Calling paper \autocite{snider2023edit} was verified through
Crossref and the institutional repository, but its own full text was not
obtained, so nothing is attributed to its findings.

The attribution is only \emph{partially} corroborated.  Petersen et al.'s
publisher-supplied abstract confirms the paper is the right kind of source
--- it reports \textcquote{petersen2015exploration}{an exploration of the EQ
parameters to determine the EQ algorithm's sensitivity to its context}
across courses taught in C, Java and Python --- but the specific values
13, 3 and 16 could not be checked against it.

Nothing \emph{refutes} the parameterisation, but the search established
something the implementation must not hide: at least four distinct EQ
parameterisations are in circulation.  Jadud's original adds 2 when both
events end in an error, 3 more for the same error type, 3 more for the same
error line and 1 more for the same edit location, and divides by 9
\autocite{jadud2005thesis}; his tuned appendix uses different constants
again; the ProgSnap2Analysis reference implementation uses 8 and 3 over 11
\autocite{price2020progsnap2}; and the two-parameter set uses 13 and 3 over
16.  A bare \enquote{EQ} number is therefore not comparable across tools,
which is why the report prints which variant produced it.

One correction to the secondary literature belongs here.  Jadud's figure is
a fall-through cascade with \emph{four} additive terms summing to a maximum
of 9, which is what makes his documented divisor of 9 correct; the widely
reproduced three-term version leaves the maximum at 8 with an unexplained
divisor.  The figure's caption mentions a maximum of 8 because it describes
a worked example whose pair happened not to share an edit location.

\paragraph{Conclusion.}
The parameterisation is real and verified; its attribution is not.  learnlog
therefore documents 13, 3 and 16 as named constants belonging to the
project's own prior work \autocite{snider2023edit}, describes them as
\enquote{following Petersen et al.}\ in the words of that work rather than
as a checked fact, and implements Jadud's original four-term scoring
\autocite{jadud2005thesis} as a documented alternative.  Obtaining the
Petersen et al.\ full text remains an open item; until then the attribution
in the prose is the manuscript's, not this project's.

\chapter{Syntax and runtime errors as separate series}
\label{app:claim:two-error-series}

\paragraph{Question.}
The Error Quotient was defined on compiler errors.  A Python run produces
both syntax errors and uncaught runtime exceptions.  Should these be scored
as one series or two?

\paragraph{Method.}
Session \texttt{learnlog-metrics-eq}, 2026-08-31.  Support query:
\qy{repeated error density compiler errors} on \texttt{dblp},
\texttt{wos}, \texttt{ieee} and \texttt{openalex}, which surfaced the
comparison study.  Counter-query: \qy{TS=(("syntactic" AND "semantic")
AND "programming behavior" AND predict*)} on \texttt{wos} ---
\textbf{empty}.  The comparison study's open preprint was read in full at
\url{https://arxiv.org/pdf/2404.05988}.

\paragraph{Results.}
One study \emph{supports} the split directly, and it did the same extension
this project needs:
\textcquote[sec.~3.4.2]{svabensky2024comparison}{Although EQ was originally
defined only for compiler errors, we extended this approach to also process
runtime errors to obtain a comparison to values for compiler errors}.  Its
result is the reason to keep the two series apart rather than merge them:
\blockcquote{svabensky2024comparison}{Compiler errors were significant
predictors of exam 1 grades but not exam 2 grades; only runtime errors
significantly predicted exam 2 grades. The findings indicate that leveraging
Error Quotient with multiple error types (compiler and runtime) may be a
better measure of students' introductory programming abilities, though still
not explaining most of the observed variability.}
The two error classes carry signal at different stages of a course;
collapsing them would discard that.

A structural argument \emph{corroborates} it from a different direction: the
Normalized Programming State Model's second axis is defined on exactly this
distinction, the presence or absence of a runtime exception in the last
execution attempt \autocite{carter2015npsm}.

Nothing \emph{refutes} the split; the dedicated counter-query was empty, and
that emptiness is recorded rather than read as endorsement.  Two
qualifications stand.  The supporting study is Java, on 280 students and
51\,095 snapshots, and its explanatory power is modest: EQ reached
\(R^2\) between 0.18 and 0.26, and homework grades outperformed all three
error measures it compared.  The one published precedent that computes these
metrics on Python data \autocite{brocker2024correlation} could not be read,
so its own syntax-versus-runtime split is unverified.

\paragraph{Conclusion.}
Two series, computed in parallel.  The evidence for the split is a single
direct study plus a structural argument, which is enough to prefer two
series over one but not enough to call it settled.  No published work argues
for restricting a Python EQ to \texttt{SyntaxError} alone, so learnlog's
choice of which exceptions count is stated in the documentation as a
documented assumption rather than as an inherited convention.

\chapter{What the Error Quotient is not}
\label{app:claim:eq-not-at-risk}

\paragraph{Question.}
Two claims frame how the output is presented.  Do these error metrics
identify at-risk students?  And do their values transfer between programming
environments, so that a student's number can be read against published
cohort figures?

\paragraph{Method.}
Session \texttt{learnlog-metrics-eq}, 2026-08-31.  This claim is where the
counter-searching was concentrated: \qy{TS=("error quotient" AND
(limitation OR criticism OR replication OR context))},
\qy{TS=("compilation behaviour" OR "compilation behavior") AND
TS=(novice)} and \qy{TS=("compilation mechanism" AND novice)} on
\texttt{wos}; \qy{error quotient context dependent limitations},
\qy{compilation error metric criticism novice programmers} and
\qy{novice programming process metrics limitations} on
\texttt{wos}, \texttt{ieee}, \texttt{dblp}, \texttt{openalex} and
\texttt{scopus}.  A separate WebSearch strand looked for any system that
shows such a metric to learners.

\paragraph{Results.}
On prediction, the strongest evidence comes from inside the same research
lineage, so it cannot be dismissed as an outsider critique.  Tabanao,
Rodrigo and Jadud built regression models from compilation behaviour and
report that \textcquote{tabanao2011predicting}{the models derived could not
accurately predict the at-risk students}.  Ahadi, Lister and Mathieson state
the general caveat: \textcquote{ahadi2018syntax}{these metrics tend to be
context dependent and contain free parameters}.  The comparison study caps
EQ's explanatory power well below the level a detector would need
\autocite{svabensky2024comparison}.

On transfer, a controlled comparison settles it.  Karvelas, Li and Becker
held population and language fixed and varied only the environment's
compilation mechanism, finding that
\textcquote{karvelas2020effects}{the programming experience and behavior of
these users can be substantially affected} by that change; the follow-up
adds that \textcquote{karvelas2022sympathy}{compilation activity differs for
each version and that students reach successful compilations faster} in the
newer environment, with a third paper corroborating from a different venue
\autocite{karvelas2020compilemuch}.  A learnlog run is an explicitly
triggered event, structurally closer to click-to-compile than to background
compilation, which supports treating a run as the compile analogue --- and
by the same result forbids comparing learnlog's EQ values with published
figures from a different environment.

The \emph{qualifying} side is real and is not suppressed.  These metrics do
carry signal: a critical review records that Jadud's work
\textcquote[p.~11:4]{hundhausen2017ide}{led to one of the first predictive
models of novice programming performance}.  The same review's proposal of a
\enquote{continuously-updated visual analytics dashboard} is its own
research agenda, not a reviewed system.  The counter-search for learner-facing
systems that show such a metric to students came back \textbf{empty}: the
closest published dashboard turns out to be a tool
\textcquote{sanchezsantillan2023dashboard}{used by the lecturer}, and
a systematic review of learner dashboards in general concludes that
\textcquote{kaliisa2024dashboards}{there is no evidence to support the
conclusion that LADs have lived up to the promise of improving academic
achievement}.

\paragraph{Conclusion.}
No to both.  The Error Quotient is not an at-risk detector and its values do
not transfer across environments, so \qy{learnlog metrics} presents a
descriptive trace of a student's own history against their own earlier
periods, never a score and never a comparison with a literature baseline.
Two caveats on the evidence: the Karvelas results are verified at abstract
strength only, so no claim is made about how much any \emph{particular}
metric shifts; and the dashboard review is general learning analytics rather
than computing-specific.  The reflective framing is what survives the
counter-search, which is also why peer comparison is not offered
\autocite{jivet2018license}.

\chapter{Repeated Error Density and its missing normaliser}
\label{app:claim:red-formula}

\paragraph{Question.}
What exactly is the Repeated Error Density formula, and is there an official
way to normalise it?

\paragraph{Method.}
Session \texttt{learnlog-metrics-eq}, 2026-08-31.  Support queries:
\qy{repeated error density compiler errors} on four providers, and
\qy{repeated error density} on \texttt{dblp}.  Counter-query:
\qy{TS=("repeated error density" AND (limitation OR criticism OR
validation OR normalis*))} on \texttt{wos} --- \textbf{empty}.  The
cross-metric query \qy{TS=("error quotient" AND ("repeated error
density" OR RED OR watwin))} on \texttt{wos} supplied the qualifying
evidence the dedicated query did not.  Becker's full text is unobtainable
(\cref{app:protocol:dropped}), so the formula was triangulated instead.

\paragraph{Results.}
Three independent sources \emph{agree exactly} on the formula.  The
open-access review states it in words:
\blockcquote{villamor2020review}{This submetric is $r_i^2/\lvert[s_i]\rvert$
where $\lvert[s_i]\rvert$ is the length of the string $s_i$ containing $r_i$
repeated errors. This can be expressed completely in terms of $r_i$ as
$r_i^2/(r_i+1)$, since the length of a string $s_i$ containing $r_i$
consecutive errors is always equal to $r_i+1$.}
That is, over the maximal strings of consecutive repeated errors in a
sequence,
\[
  \mathrm{RED} = \sum_{i=1}^{n} \frac{r_i^{2}}{\lvert s_i\rvert}
               = \sum_{i=1}^{n} \frac{r_i^{2}}{r_i + 1}.
\]
The ProgSnap2Analysis reference implementation computes exactly this, and
the worked values in the comparison study discussed below match.  Becker's
own abstract supplies the motivation, verified through two providers: RED
\textcquote{becker2016new}{has advantages over EQ including being less
context dependent, and being useful for short sessions}.

On normalisation the answer is negative and well attested.  The comparison
study's summary table gives RED's range as \([0,\infty)\) against EQ's
\([0,1]\) \autocite{svabensky2024comparison}; and the reference
implementation says in its own source that there is
\enquote{no official way to normalize RED}, dividing by the number of
compiles as its own choice --- a decision of the reimplementation, not of
the paper.

The counter-evidence \emph{qualifies} RED sharply.  The review's comparison
table lists it as unbounded, outlier-prone and not yet validated against
student performance \autocite{villamor2020review}, and the same comparison
study finds RED predicting course grades \emph{worse} than EQ.

\paragraph{Conclusion.}
The formula is established --- by triangulation across three sources that
agree exactly, not by reading the original --- and there is no official
normaliser.  learnlog therefore reports RED unnormalised and prints the run
count beside it, so the number can be read at all, and does not present RED
as an improvement on EQ.  Its reason for being included is narrower and
better supported: Jadud required at least seven events for a session, and
RED was designed for short sessions, which is what a single project's run
history usually offers.

\chapter{Cyclomatic complexity: predicates plus one}
\label{app:claim:cc-definition}

\paragraph{Question.}
learnlog counts a function's cyclomatic complexity as decision points plus
one, adding one per condition in a compound boolean.  Are both rules the
primary source's own, or are they the measuring tool's convention?

\paragraph{Method.}
Session \texttt{learnlog-metrics-code}, 2026-08-31.  Support query:
\qy{cyclomatic complexity measure control flow graph} on \texttt{dblp},
\texttt{wos}, \texttt{ieee} and \texttt{openalex}.  McCabe's own paper did
not appear in the ranked hits --- a good illustration of relevance ranking
burying an older known item --- so it was retrieved by DOI through Crossref
and its full text obtained from an open copy at
\url{http://www.literateprogramming.com/mccabe.pdf} and read (13 pp.).
Counter-query: \qy{TS=("cyclomatic complexity" AND (critique OR
criticism OR "empirical validation" OR invalid))} on \texttt{wos}.  The
natural-language form of that counter-query,
\qy{cyclomatic complexity critique validity}, returned zero hits and is
recorded as a failed query rather than as a negative result.

\paragraph{Results.}
Both rules are McCabe's own words.  The measure is defined graph
theoretically --- \textcquote[p.~308]{McCabe1976}{The cyclomatic number V(G)
of a graph G with n vertices, e edges, and p connected components is v(G) =
e - n + p} --- and then reduced to the form an implementation can use:
\textcquote[p.~315]{McCabe1976}{the cyclomatic complexity of a structured
program equals the number of predicates plus one}.  The compound-predicate
rule is stated just as explicitly:
\blockcquote[p.~315]{McCabe1976}{In practice compound predicates such as IF
\enquote{c1 AND c2} THEN are treated as contributing two to complexity since
without the connective AND we would have IF c1 THEN IF c2 THEN which has two
predicates.}

The counter-search \emph{did not} touch the definition.  It surfaced
Shepperd's critique \autocite{Shepperd1988} and the empirical re-tests
\autocite{Landman2016}, all of which dispute the measure's validity as a
proxy for complexity or maintainability --- the subject of
\cref{app:claim:cc-sloc} --- while taking the definition as given.

\paragraph{Conclusion.}
Yes: both counting rules come from the primary source, so the documentation
cites McCabe and not the measuring tool for them.  The caveat is a boundary,
not a doubt: what is uncontested is the \emph{definition}.  What the same
counter-search shows to be contested is what the number means, which the
next chapter takes up.  learnlog's implementation additionally documents the
tool-specific increments --- for loop and try clauses, comprehensions and
match statements --- from the tool's source rather than from its
documentation, because the two disagree.

\chapter{Is cyclomatic complexity redundant with size?}
\label{app:claim:cc-sloc}

\paragraph{Question.}
It is often said that cyclomatic complexity carries no information beyond
lines of code.  Is that established?

\paragraph{Method.}
Session \texttt{learnlog-metrics-code}, 2026-08-31.  Both sides came from
two field-tagged Web~of~Science queries: \qy{TS=("cyclomatic complexity"
AND "lines of code" AND correlat*)} and \qy{TS=("cyclomatic complexity"
AND (critique OR criticism OR "empirical validation" OR invalid))}.  This is
a two-sided question by construction, so the same query was read as both
support and counter-search, and both positions are cited.

\paragraph{Results.}
The redundancy position is stated most strongly by Jay et al., whose models
\textcquote{JayEtAl2009}{successfully predict roughly 90\% of CC's variance
by LOC alone}, concluding that the measure has \enquote{absolutely no
explanatory power of its own}.  Read at abstract strength; the venue is a
weak one, which is a further reason not to rest a claim on it alone.

That position is \emph{contradicted} directly, and at much larger scale, by
Landman et al., on 17.6 million Java methods and 6.3 million C functions:
\blockcquote{Landman2016}{Our results show that linear correlation between
SLOC and CC is only moderate as a result of increasingly high variance. We
further observe that aggregating CC and SLOC as well as performing a power
transform improves the correlation. Our conclusion is that the observed
linear correlation between CC and SLOC of Java methods or C functions is not
strong enough to conclude that CC is redundant with SLOC. This conclusion
contradicts earlier claims from literature.}
The disagreement is not a tie: it turns on \emph{aggregation}.  The
correlation looks strong when methods are aggregated and weak at the level
of the individual method --- which is the level at which learnlog reports.

Two further findings \emph{qualify} the measure without settling the
redundancy question.  Shepperd's theoretical objections stand: he attacks
\textcquote{Shepperd1988}{the very simplistic approach to decision counting},
since the ease of comprehending a decision \enquote{is not invariant}, and
\enquote{the arbitrary impact of modularisation upon total program
complexity}.  And at novice scale the measure is
near-degenerate: as reported by Keuning et al., in
\textcquote[p.~2]{Keuning2017}{78\% of over 4 million scripts the cyclomatic
complexity is one. Only 4\% of the scripts has a complexity over four.}

\paragraph{Conclusion.}
No, it is not established --- the redundancy claim is contested, and the
larger and more recent study contradicts it at exactly the granularity
learnlog uses.  The documentation therefore states the disagreement rather
than picking a side, reports complexity \emph{with} source lines and never
alone, and reports the per-function maximum and the count above one rather
than a mean.  Caveats: neither study concerns novice code or Python, and the
78\,\% figure is from Scratch and is cited as reported by Keuning et al.,
the primary being closed access.  The practical consequence for a beginner's
project is that complexity will sit at one for weeks; the informative event
is when it first rises.

\chapter{The maintainability index needs its caveats}
\label{app:claim:mi-caveats}

\paragraph{Question.}
Should a maintainability index be reported to a student, and if it is
reported at all, what must accompany it?

\paragraph{Method.}
Session \texttt{learnlog-metrics-code}, 2026-08-31.  Support query:
\qy{maintainability index software metrics} on \texttt{dblp},
\texttt{wos}, \texttt{ieee} and \texttt{openalex}.  The natural-language
counter-query \qy{maintainability index criticism limitations} returned
only off-topic records and is recorded as a failed query; the critical
literature came from the same support query and from the field-tagged
complexity queries.  Coleman et al.'s full text was retrieved and page~46
read directly.

\paragraph{Results.}
The formula's own authors published the first caveat.  Having given the
polynomial --- 171 less 5.2 times the natural log of average volume, less
0.23 times average complexity, less 16.2 times the natural log of average
lines, plus a comment term --- they write:
\blockcquote[p.~46]{ColemanEtAl1994}{Preliminary results indicated that this
model was too sensitive to large numbers of comments. That is, large comment
blocks, especially in small modules, unduly inflated the resulting
maintainability indices.}
Small, heavily commented modules are exactly what a beginner writes.

Heitlager et al.\ attack the constants and the comment term:
\textcquote{Heitlager2007}{there is no logical argument why the MI formula
contains the particular constants}, asking why complexity is multiplied by
0.23, and concluding that
\textcquote{Heitlager2007}{counting the number of lines of comment, in
general, has no relation with maintainability whatsoever}.  Sjøberg et al.\
tested such surrogate measures against \emph{measured} maintenance effort on
four functionally equivalent systems and found that
\textcquote{Sjoberg2012}{the metrics were not mutually consistent} and that
\textcquote{Sjoberg2012}{apart from size, surrogate maintainability measures
may not reflect future maintenance effort}.  Morrison et al.\ ask the
question closest to this project's --- the index's accuracy for Python as a
function of module size \autocite{Morrison2023} --- but the paper is closed
access, so only the existence of the question is verified, not the direction
of the effect.

No source \emph{defending} the index against measured effort was found; the
counter-search for one came back empty.  Its adoption in mainstream
development tools shows implementation, not validity.  The origin of the
index \autocite{OmanHagemeister1992} is recorded as an attribution only, its
content unread.

\paragraph{Conclusion.}
Report it only with its caveats attached, and never as a headline.  The
supporting evidence is one-sided and includes the original authors' own
warning.  The honest scope limit is that none of these studies examined
novice or tiny programs: the argument that the index is meaningless for a
fifteen-line script rests on Coleman et al.'s small-module admission, on
Morrison et al.\ having thought the question worth asking, and on the
arithmetic of the implementation, which returns a perfect score when
Halstead volume or source lines are zero --- as they are for a first
program.  No study of novices was found either way.  The caveat is therefore
printed in the report header rather than left to the documentation.

\chapter{Linter rule sets for novices}
\label{app:claim:linter-rules}

\paragraph{Question.}
Can a general-purpose linter's default rule set be pointed at student code,
or must the rules be selected?

\paragraph{Method.}
Session \texttt{learnlog-metrics-code}, 2026-08-31.  Support queries:
\qy{code quality student programs} on \texttt{dblp}, \texttt{wos},
\texttt{ieee} and \texttt{openalex}, which returned the decisive study as
its first hit; and \qy{static analysis introductory programming
feedback} on the same providers, which doubled as the counter-search by
looking for evidence that student-facing static analysis \emph{does} work.
The decisive study's full text was read from the authors' institutional
technical report.

\paragraph{Results.}
The claim is \emph{supported} by the largest study of the question, over
453\,526 student files and 2.66 million snapshots, whose own recommendation
is the design rule:
\blockcquote{Keuning2017}{many of the checks PMD can perform are not
suitable for novice programmers \ldots{} We advise educators to customize the
tool by selecting a small set of relevant issues.}
Its supporting figures make the point concrete: the most frequent raw
violations are pedagogically irrelevant, students rarely fix what is
reported, and students using a code-analysis extension had slightly
\emph{more} issues per thousand lines than those using none.

An independent Python precedent \emph{corroborates} the restrict-and-rewrite
pattern: PyTA wraps a general-purpose linter with custom checks and improved
messages for novices \autocite{LiuPetersen2019}.  Its authors' own scope
limit is preserved here --- the improvement was observed only for students
who chose to read the output, a self-selected subgroup, so no causal claim
is transferred.  Stegeman et al.\ supply a vocabulary for grouping what is
reported, in three families covering documentation, presentation and
structure \autocite{Stegeman2016}; that category list was read from Keuning
et al.'s restatement, not from the rubric paper itself.

The counter-search found positive results for student-facing static analysis
in general, which narrows the claim rather than refuting it: what is
unsuitable is the \emph{default rule set}, not the practice of showing
diagnostics.  It also found the other direction, in a deployed classroom
study reporting that
\textcquote{Bobadilla2024}{students add more violations than they remove}.

\paragraph{Conclusion.}
The rules must be selected.  learnlog therefore runs a small, named rule set
aimed at mistakes a beginner can act on --- undefined names, unused imports
and variables, bare excepts, comparisons to singletons --- rather than a
tool's defaults, and treats the linter as optional, reporting its absence
instead of failing.  Caveats: the decisive study is Java, one analyser and
one dataset; the rubric's own text was not read; and the evidence that
quality feedback changes novice behaviour is genuinely mixed, which is why
the output is descriptive rather than corrective.

\chapter{Headline aggregates mask the long tail}
\label{app:claim:long-tail}

\paragraph{Question.}
Is a single improving summary number an adequate way to report code quality
over time to a student?

\paragraph{Method.}
Session \texttt{learnlog-metrics-code}, 2026-08-31.  Query:
\qy{TS=("code quality" AND (students OR novice*) AND (evolution OR "over
time" OR trend* OR longitudinal))} on \texttt{wos}.  A related field-tagged
query, \qy{TS=("software metrics" AND (students OR novice OR
"introductory programming"))}, was run because an earlier pre-protocol pass
using natural-language queries without Web~of~Science had concluded that no
literature existed on metric trends for students.  That conclusion was a
\textbf{false negative}, and the query overturned it --- which is the
clearest evidence in this appendix for why the provider and syntax choices
in \cref{app:protocol:sessions} matter.

\paragraph{Results.}
The claim is \emph{supported} directly by the longest-running longitudinal
measurement found, over four course iterations, 200 students and 19
assignments:
\blockcquote{Ostlund2023}{Results showed a statistically significant
improvement in code quality over 19 assignments in all four course
iterations. In particular, when a single violation was the focus of a
learning outcome, the effects were both a dramatic fall and sustained low
occurrence. However, this improvement was mostly focused on the three most
frequently occurring violations ($>$~70\%), which masked an increase in lesser
occurring violations throughout the course.}
The headline improved while part of the distribution got worse.

Two studies \emph{corroborate} the pattern from other angles.  Prajapati and
Acuña, tracking quality per module across a course, find improvement overall
while one category stays dominant:
\textcquote{Prajapati2025}{Formatting errors, such as whitespace violations,
remain high, comprising nearly half of all detected issues}.  Bobadilla et
al.\ find net violations rising \autocite{Bobadilla2024}.  Neither reports to
students --- both address researchers or instructors --- so they establish
the measurement, not the presentation.

On presentation the evidence is thin and honestly labelled as such.  Van den
Aker and Rahimi derive nine design principles for generating and presenting
automated formative feedback on code quality,
\textcquote{vandenAker2024}{rooted in solid educational constructs about
feedback and learning dashboards}; but the full text was unobtainable, so
this backs
the existence of such principles and nothing about their content.  The gap
that survives the corrected search is specific: metric \emph{trends} shown
to students.

\paragraph{Conclusion.}
No --- a lone headline aggregate is not adequate, and the one longitudinal
study of exactly this question shows why.  \qy{learnlog metrics}
therefore reports distributions and long-tail counts beside any summary, and
describes direction in neutral words without evaluating the student.
Caveats: the masking finding is a single course sequence in a single
institution, and the design guidance for presenting such trends to students
is cited unread and remains to be obtained.

% Restore url.sty's default: no line breaks after hyphens.
\makeatletter
\def\do@url@hyp{}
\makeatother
\Urlmuskip=0mu\relax
